\documentclass[12pt,a4paper]{article}  % Déclare la classe du document.

%  ------------------------------------------  Modules usuels  --------------------------------------
%  --------------------------------------------------------------------------------------------------------

\usepackage{amsmath} 	
\usepackage{amssymb} 
\usepackage{listings}         
\usepackage{numprint} 
\usepackage{pifont}

%  --------------------------------------------------------------------------------------------------------
% ----------------------  packages pour fontes et caractères spéciaux  ---------------------
%  --------------------------------------------------------------------------------------------------------

\usepackage{eurosym}   
\usepackage{amsfonts}   
\usepackage{amstext}       
\usepackage{amsfonts}	
\usepackage{mathrsfs}	 

\usepackage{pstricks,pst-all,pstricks-add,pst-tree,pst-node,pst-3dplot}

%------------------------------------------------------------------------------------------------------
%  --------------------------  Francisation et gestion des caractères sous mac --------------------------------------------------------------------------------------------------------

\usepackage[francais]{babel}   
\usepackage[utf8]{inputenc}
\usepackage[OT1,T1]{fontenc} 
%  --------------------------------------------------------------------------------------------------------
% ------------------------------------Packages  de  mise en page ---------------------------
%--------------------------------------------------------------------------------------------------------

\usepackage{fancyhdr}	% en-têtes et pieds de pages
\usepackage{amsfonts}

 \usepackage{multicol} % permet l'utilisation du multicolonnage dans le texte. La syntaxe est \begin{multicols}{nb de colonnes}    ......  \end{multicols}
 
\usepackage{calc}            % permet de faire un minimum de calcul sur les arguments des commandes et macros (utilisé dans la def de l'environnement preuve )
\usepackage{epic}            % extension et simplification de  l'environnement picture
\usepackage{graphicx}    % permet d'inclure des graphique externe dans un document
\usepackage{variations}   % permet d'obtenir de joli tableau de variations
\usepackage{enumerate}  %permet de faire des listes personnalisées
\usepackage{arydshln}     %permet d'avoir des lignes en pointillés
\usepackage{framed}  
\usepackage{multido} 


%  --------------------------------------  MACROS USUELLES  ----------------------------------------

\newcommand{\R}{\mathbb{R}} 
\newcommand{\N}{\mathbb{N}} 
\newcommand{\D}{\mathbb{D}} 
\newcommand{\Z}{\mathbb{Z}} 
\newcommand{\Q}{\mathbb{Q}} 
\newcommand{\C}{\mathbb{C}} 

%  -------------------------------------  MACROS SPECIFIQUES ------------------------------------


\newcommand{\raisedrule}[2][0em]{\leaders\hbox{\rule[#1]{1pt}{#2}}\hfill}


\newcounter{numexo}

\setcounter{numexo}{1}

\newcommand{\exercice}[1]{\vskip 0.6cm\noindent\rule[+2pt]{1cm}{0.5pt}\ {\large\bf \ Exercice \thenumexo {  #1 }}\raisedrule[+2pt]{0.5pt}\null\addtocounter{numexo}{1}\par\vskip 0,4cm}


\newcommand{\lignes}[1]{
\setlength{\unitlength}{1mm}
\vskip 0,4cm\noindent
\multido{\i=0+1}{#1}{%
     \noindent
     \begin{picture}(150,5)
     \linethickness{0,005mm}
     \put(-5,0){\dashline{1}(175,5)(5,5)}
     \end{picture}\vskip 0,001cm
     }
\vskip 0cm
}




% ----------------------------------------- FORMAT DE PAGE  -----------------------------------------

\usepackage[left = 1.3cm,
	 	      right = 1.3cm,
		      top = 1.7cm,
		      headsep = 8mm,
		      bottom = 1.7cm,
		      headheight = 14.2 pt
		      ]{geometry}

\setlength{\parindent}{0mm}


\everymath{\displaystyle} 
\renewcommand{\thesection}{\Roman{section}} 
\author{}
\title{}
\date{}
\pagestyle{fancy}

 
\lhead{1ère - spécialité maths} 


\rhead{DS n$^{\circ}$7 - 18/05/21}		


\fancyfoot[C]{} 
\fancyfoot[L]{} 
\fancyfoot[R]{} 

\lstset{language=Python}  

\lstset{frame = single
framerule = 1pt
showspaces = false,
showstringspaces = false,
numbers=left, numberstyle=\tiny}

% ----------------------------------------- correction  -----------------------------------------
\usepackage[most,skins,breakable]{tcolorbox}
\usepackage{enumitem}

\newcommand{\corr}[1]{\fcolorbox{red}{yellow!30}{#1}}

\newcommand{\bareme}[1]{\tikz\node[circle,draw,inner sep=1pt,outer sep=1pt, color=red]{\textbf{#1}};}

\newtcolorbox{corrBox}{enhanced,boxrule=1pt,nobeforeafter,boxsep=0pt,left=\fboxsep,right=\fboxsep,colback=yellow!30,colframe=red!75!black}

\newlist{enumMet}{enumerate}{2}
\setlist[enumMet,1]{label=\arabic*), font=\bf\color{blue!75!black}}
\setlist[enumMet,2]{label=\alph*), font=\bf\color{blue!75!black}}


\begin{document}
\graphicspath{{images/}}

\begin{center}

\fbox{\LARGE{\rule{0.5cm}{0pt}\rule[-0.25cm]{0pt}{0.9cm} Devoir Surveillé N$^{\circ}$7 - Application de la dérivation	\rule{0.5cm}{0pt}}}
\end{center}


\begin{center}\textbf{1h - Calculatrice autorisée - Barème indicatif}\end{center}

\begin{center}\textbf{Les élèves avec un tier-temps ne traitent pas les questions avec  le symbole \ding{46}}\end{center}

\begin{center}
%On veillera à ENCADRER ses résultats et soigner sa copie.\\
Les résultats doivent être justifiés par des calculs (au moins 1 étape intermédiaire).
\end{center}
\vspace{.05cm} 

\vspace{-0.3cm}
\exercice{ - Tableau de variations - ( 10 points )}

Pour chaque fonction, définie et dérivable sur I, Calculer sa dérivée, puis dresser son tableau de variations
\begin{enumMet}
    \item $f: x \mapsto 3x^{2}-6 x+9, \mathrm{I}=\mathbb{R}$
    \item $g: x \mapsto x^{3}-6 x^{2}+9 x-1, \mathrm{I}=\mathbb{R}$
    \item $h: x \mapsto \frac{3 x-1}{x+2}, I=\mathbb{R} \backslash\{-2\}$.
    \item \ding{46} $k: x \mapsto x^{4}-\frac{4}{3} x^{3}, \mathrm{I}=\mathbb{R}$
\end{enumMet}

\exercice{ - extremums locaux - ( 5 points )}
On suppose que $f$ et $g$ sont des fonctions définies et dérivables sur un intervalle donné. Déterminer les éventuels extremums locaux des fonctions $f$ et $g$ sur cet intervalle.
\begin{enumMet}
    \item $f(x)=2x(3x^2-18)$
    \item $g(x)=4x-5+\dfrac{3}{x-2}$
\end{enumMet}

\exercice{ - Recette et bénéfice - ( 6 points )}

Une entreprise fabrique et vend des trottinettes électriques. On modélise le coût total de fabrication journalier, en centaines d'euros, de $x$ trottinettes ($2\leqslant x\leqslant 14)$ par la fonction notée $C$ telle que :
$$
C(x)=\frac{1}{3} x^{3}-6 x^{2}+36 x-50
$$
Le nombre $x$ de trottinettes demandées est fonction de leur prix unitaire $P$, qui vérifie  la condition $P=12-\frac{1}{2} x$.

\begin{enumMet}
    \item \begin{enumMet}
        \item Déterminer la recette totale journalière $R(x)$ correspondant à la vente de $x$ trottinettes.
        \item Dresser le tableau de variations de la fonction $R(x)$
        \item Pour quel nombre de trottinettes vendues quotidiennement réalise-t-on la recette maximale ? Quel est dans ce cas le prix d'une trottinette?
    \end{enumMet}
    \item \begin{enumMet}
        \item Exprimer le bénéfice journalier $B(x)$ obtenu pour la vente de $x$ trottinettes.\\
        \emph{Rappel : Le bénéfice d'une entreprise se calcule par différence entre la recette et le cout de production.}
        \item \ding{46} Dresser le tableau de variations de la fonction $B(x)$
        \item \ding{46} Pour quel nombre de trottinettes vendues réalise-t-on le bénéfice journalier maximal ? Quel est dans ce cas le prix d'une trottinette électrique?
    \end{enumMet}
    % \item \textbf{Question Bonus}
    % \begin{enumMet}
    %     \item Avec la calculatrice, représenter dans un même repère les courbes représentatives des fonctions $C$, $R$ et $B$ sur [2; 14].
    %     \item Que peut-on dire du bénéfice journalier de cette entre prise lorsque la recette journalière est maximale? Justifier.
    % \end{enumMet}
\end{enumMet}

\end{document}

